\documentclass[12pt,a4paper]{article}

\usepackage[utf8]{inputenc}
\usepackage[T2A]{fontenc}
\usepackage[serbianc]{babel}

\usepackage{amsmath,amsthm}
\usepackage{graphicx}
\usepackage{wrapfig}
\usepackage{xcolor}
\usepackage{geometry}
\usepackage{array}
\usepackage{ulem}
\usepackage{xcolor}

\usepackage[hidelinks]{hyperref}



\geometry{margin=2.5cm}

\definecolor{myblue}{HTML}{659EC7}

\newtheorem{definition}{Дефиниција}
\newtheorem{theorem}{Теорема}
\newtheorem{lemma}{Лема}

\begin{document}

\begin{titlepage}

\centering

\vspace*{2cm}

{\large\bfseries УНИВЕРЗИТЕТ У БЕОГРАДУ}\\[0.3cm]
{\large\bfseries МАТЕМАТИЧКИ ФАКУЛТЕТ}

\vspace{2cm}

{\large СЕМИНАРСКИ РАД}\\[0.3cm]
{\large из предмета}
{\Large\bfseries Увод у информатику}

\vspace{2.5cm}
\rule{0.75\textwidth}{0.8pt}

\vspace{0.6cm}

{\Large\bfseries МАГИЧНИ КВАДРАТИ}\\[0.35cm]

{\Large\itshape Примена математичких принципа}\\[0.15cm]
{\Large\itshape у рачунарству}

\vspace{0.6cm}

\rule{0.75\textwidth}{0.8pt}
\vspace{5cm}

\begin{flushleft}
\begin{tabular}{ll}
\textbf{Ментор:} & Данијела Симић \\[0.25cm]
\textbf{Студент:} & Владана Милетић\\[0.25cm]
\textbf{Број индекса:} & 213/2025
\end{tabular}
\end{flushleft}

\vspace{1.5cm}

{\large Београд, jул 2026.}

\end{titlepage}



\tableofcontents
\newpage



\section{Увод}


\textbf{Магични квадрати} представљају занимљиве математичке структуре које
повезују бројеве, симетрију и логику. Они су један од најпознатијих примера
дискретне математике.


\textcolor{myblue}{Магични квадрати показују како једноставна правила могу створити
сложене математичке структуре.}


Циљ овог рада је да прикаже основне појмове, особине и примене магичних
квадрата.



\section{Основни појмови}



\subsection{Дефиниција магичног квадрата}



\begin{definition}

Магични квадрат реда $n$ је квадратна табела димензије $n \times n$ у којој је
збир бројева у сваком реду, колони и дијагонали једнак. Та вредност назива се
магична константа.

\end{definition}



Магична константа се израчунава формулом:


\[
M=\frac{n(n^2+1)}{2}
\]


За магични квадрат трећег реда:


\[
M=\frac{3(3^2+1)}{2}=15
\]

\subsection{\text{Магични квадрат 3 × 3}}

{\color[HTML]{659EC7}
Најпознатији магични квадрат реда \(3 \times 3\) састоји се од бројева од 1 до 9, при чему се сваки број појављује тачно једном. Збир бројева у сваком реду, свакој колони и на обе дијагонале увек је једнак \textbf{15}. Тај број назива се \textbf{магични збир}.
}

\vspace{0.5cm}

\begin{center}
{\color[HTML]{659EC7}
\renewcommand{\arraystretch}{2}
\setlength{\tabcolsep}{18pt}
\begin{tabular}{|c|c|c|}
\hline
\Large\textbf{8} & \Large\textbf{1} & \Large\textbf{6} \\ \hline
\Large\textbf{3} & \Large\textbf{5} & \Large\textbf{7} \\ \hline
\Large\textbf{4} & \Large\textbf{9} & \Large\textbf{2} \\ \hline
\end{tabular}
}
\end{center}

\vspace{0.3cm}

\begin{center}
{\large\bfseries\color[HTML]{659EC7}Магични збир = 15}
\end{center}








\section{Математичке особине}



\begin{lemma}

У нормалном магичном квадрату реда $n$ бројеви од $1$ до $n^2$ појављују се
тачно једанпут.

\end{lemma}



\begin{theorem}

Збир свих елемената магичног квадрата реда $n$ једнак је:

\[
S=nM
\]

где је $M$ магична константа.

\end{theorem}




\section{Примене магичних квадрата}



Магични квадрати имају примену у различитим областима:


\begin{itemize}

\item дискретној математици,

\item комбинаторици,

\item развоју алгоритама.

\end{itemize}



Према начину конструкције могу се поделити на:


\begin{enumerate}

\item магичне квадрате непарног реда,

\item магичне квадрате парног реда,

\item посебне симетричне квадрате.

\end{enumerate}



\begin{description}

\item[Нормални магични квадрат]
садржи бројеве од $1$ до $n^2$.

\item[Полумагични квадрат]
има једнаке збирове само у редовима и колонама.

\end{description}





\section{Приказ магичног квадрата}

\flushleft
\vspace{1cm}

\includegraphics[width=0.6\textwidth]{magic3.png}

\vspace{1cm}

Магични квадрати се могу приказати и графички, што омогућава лакше
разумевање њихове симетрије и правилности.



\newpage


\section{Закључак}



Магични квадрати представљају занимљив спој математике, логике и симетрије.
Иако су настали као математичке загонетке, данас имају примену у различитим
областима рачунарства и математике.
\vspace{0.2cm}

\textcolor{red}{\uline{Ова тема показује да се из једноставних математичких правила могу добити веома сложене и занимљиве структуре}}.


\end{document}